% !PRINTED 266
\weberpage{266}
%<<CONT
would have to be, and $a$ and $d$ would not come out real. Accordingly
we take for $a$ and $c$, as above, the $p\,(p-1)$ different
systems of values, and obtain from (11) for each of them
$1/2\,(p-1)$ determinations of $b$ and $d$; this number too is to be
halved, and there result
\[
1/4\;p\,(p-1)\,(p-1)
\]
substitutions of the III\textsuperscript{rd} kind. The total number of all substitutions
 of the group $L_p$ is accordingly, as it must be,
\[
p^2 + 1/4\,p\,(p-3)\,(p+1) + 1/4\,p\,(p-1)^2 = 1/2\,p\,(p^2-1).
\]

\begin{center}
§. 69.

\gesperrt{Imaginary form of the group $L_p$.}
\end{center}

The substitutions of the first two kinds of the group $L_p$
have the property that the normal form $T$ or $S$ corresponding
 to them is likewise contained in $L_p$, and that one can
transform them into these normal forms by \gesperrt{real} substitutions
, that is, by substitutions which themselves occur in $L_p$
. Neither of these holds for the substitutions of the third
kind.

One can, however, as we shall now prove, transform the whole
group $L_p$ into another group $\Lambda_p$ in such a way that $\Lambda_p$ is
\gesperrt{a divisor of the group $E_{2,p}$ isomorphic with $L_p$}, that the
normal forms of the transformations of the third kind are contained
in $\Lambda_p$, and that every substitution of the third kind in $\Lambda_p$ can
be transformed into the normal form by substitutions of
$\Lambda_p$ itself.

In order to find such a transformation of $L_p$, we consider
the substitution
\weq{1}{S = \begin{pmatrix} \sigma, & 0 \\ 0, & \sigma^{-1} \end{pmatrix}}
under the assumption that $\sigma$ and $\sigma^{-1}$ are conjugate imaginary
, and seek to determine the substitution
\weq{2}{R = \begin{pmatrix} \lambda, & \mu \\ \nu, & \varrho \end{pmatrix}}
in such a way that
\weq{3}{R\,S\,R^{-1}}
%CONT>>
